\subsection{A finite obstruction to a derived-binomial comparison}

The associated graded of the derived binomial monad is the sum of
nonabelian derived divided powers. In the current arXiv version of
Kubrak--Shuklin--Zakharov \cite{A_KSZ}, this is Remark~2.52; Lemma~2.54 and
Remark~2.55 describe computation of each finite filtration stage by
finite cosimplicial totalization. This observation suggests applying
cosimplicial denormalization to the two-term complex
$[\Phi\to\Psi]$. We now test the resulting divided-power complex
before taking derived sections. The following obstruction concerns
that specific construction. It does not contradict the derived
binomial theorem, and does not exclude every possible filtered
binomial model.

\begin{theorem}[A-022: \status{Refuted}; the naive cosimplicial bridge]
\label{sa08:A022}
Let $k=\Z$, let $\Phi_n=\Z^{n+1}$ be the standard simplicial module
with marked constant vector $T$, and put $\Psi=\Phi/\Z T$.
Let $M=[\Phi\xrightarrow{d}\Psi]$, in cochain degrees zero and one,
and define complexes of simplicial abelian groups
\[
 K=\bigl[\Gamma^2\Phi\longrightarrow\Phi\otimes\Psi
                  \longrightarrow\Lambda^2\Psi\bigr],\qquad
 C=N_{\mathrm{cosimp}}\Gamma^2
                 (\operatorname{DK}_{\mathrm{cosimp}}M).
\]
Here $K$ is the weight-two standard PD de Rham complex. There is a
natural, simplexwise quasi-isomorphism $\theta:C\to K$, but there
is no cochain homotopy equivalence between $C$ and $K$ in
$\operatorname{Ch}(s\operatorname{Mod}_{\Z})$.
More strongly, for the finite simplicial set
$X=\operatorname{sk}_2 B(\Z/2)$ one has
\[
 H^2\operatorname{Hom}(X,C)=0,\qquad
 H^2\operatorname{Hom}(X,K)=\Z/2.
\]
Thus replacing the PD Koszul complex by this objectwise
cosimplicial divided-power model does not preserve underived
sections on every simplicial set. For
$X=\Delta[2]/\partial\Delta[2]$, the map induced by $\theta$ on
second cohomology is multiplication by two between infinite cyclic
groups.
\end{theorem}

\begin{proof}
We first calculate $C$ without a comparison theorem. In cosimplicial
degree $r$ the denormalization of $M$ is
$\Phi\oplus\Psi^{\oplus r}$. Its first coface sends
$(a,b_1,\ldots,b_r)$ to $(a,da,b_1,\ldots,b_r)$, its internal
cofaces repeat one $b_i$, and its last coface appends zero.
A codegeneracy omits one $\Psi$ coordinate. Consequently the
normalized second divided power has precisely the terms
\[
 C^0=\Gamma^2\Phi,\qquad
 C^1=(\Phi\otimes\Psi)\oplus\Gamma^2\Psi,\qquad
 C^2=\Psi\otimes\Psi.
\]
There are no further terms: a quadratic monomial cannot involve
three distinct cosimplicial coordinates. If
$\delta:\Gamma^2\Psi\to\Psi\otimes\Psi$ denotes the
$(1,1)$-part of the divided-power diagonal, the differentials are
\begin{align*}
 d_C\gamma_2(a)&=(a\otimes da,\gamma_2(da)),\\
 d_C(ab)&=(a\otimes db+b\otimes da,da\,db),\\
 d_C(a\otimes b,c)&=da\otimes b-\delta(c).
\end{align*}
The Koszul differentials are the first components of the first two
formulas and $d_K(a\otimes b)=da\wedge b$.
Thus $\theta$ is the identity in degree zero, projection onto
$\Phi\otimes\Psi$ in degree one, and the exterior quotient in
degree two. Since every $\Psi_n$ is free, the sequence
\[
 0\longrightarrow\Gamma^2\Psi
 \xrightarrow{\delta}\Psi\otimes\Psi
 \longrightarrow\Lambda^2\Psi\longrightarrow0
\]
is exact. Indeed, in a free basis $e_i$, the first map sends
$\gamma_2(e_i)$ to $e_i\otimes e_i$ and $e_ie_j$ to
$e_i\otimes e_j+e_j\otimes e_i$ for $i<j$.
The kernel complex of $\theta$ is therefore the identity on
$\Gamma^2\Psi$, with a minus sign, in degrees one and two.
This proves that $\theta$ is a simplexwise quasi-isomorphism.
It does not prove a homotopy equivalence in simplicial modules:
the displayed short exact sequence need not split there.

Here are full integral normalization matrices that detect the
failure. Write
\[
 \Phi_n=\Z T\oplus\bigoplus_{j=1}^n\Z u_j,
 \qquad \Psi_n=\bigoplus_{j=1}^n\Z u_j,
 \qquad d(T)=0,\quad d(u_j)=u_j.
\]
We use the same letter for the image of $u_j$ in $\Psi_n$.
On $\Phi$, the first face sends $u_1$ to $T$ and $u_j$ to
$u_{j-1}$ for $j>1$; on $\Psi$ it sends $u_1$ to zero instead.
An internal face identifies $u_i,u_{i+1}$, and the last face sends
$u_n$ to zero. Every face fixes $T$.
Normalize simplicially by $N_sY_n=\bigcap_{i=1}^n\ker(d_i)$,
with differential $d_0$. For compactness, use the expressions
\begin{align*}
 b&=Tu_1, & c&=\gamma_2(u_1),
 &h&=u_1u_2-2\gamma_2(u_2),\\
 p&=T\otimes u_1, & q&=u_1\otimes u_1,
 &\eta&=u_1\wedge u_2,\\
 r&=u_1\otimes u_2-u_2\otimes u_2,
 &s&=u_2\otimes u_1-u_2\otimes u_2.
\end{align*}
An expression is interpreted in the ambient module in its row.
The following bases give all nonzero normalized terms:
\[
\begin{array}{c|c|c|c}
Y&N_sY_0&N_sY_1&N_sY_2\\ \hline
\Gamma^2\Phi&\gamma_2(T)&b,c&h\\
\Phi\otimes\Psi&0&p,q&r,s\\
\Gamma^2\Psi&0&c&h\\
\Psi\otimes\Psi&0&q&r,s\\
\Lambda^2\Psi&0&0&\eta
\end{array}
\]
These bases can be verified directly from $d_1,d_2$.
There are no normalized terms in degrees greater than two.
Indeed a simplicial degeneracy omits one of the $u_j$ coordinates;
each quadratic monomial involves at most two such coordinates, so
in higher dimensions it is degenerate. The quotient by degenerates
is canonically isomorphic to the displayed normalization.
In the order of the table, the nonzero simplicial differentials are
\[
\begin{array}{c|c|c}
Y&N_sY_1\longrightarrow N_sY_0
 &N_sY_2\longrightarrow N_sY_1\\ \hline
\Gamma^2\Phi&(\,2\quad1\,)&\binom{1}{-2}\\
\Phi\otimes\Psi&0&
 \begin{pmatrix}1&0\\-1&-1\end{pmatrix}\\
\Gamma^2\Psi&0&(-2)\\
\Psi\otimes\Psi&0&(\,-1\quad-1\,)\\
\Lambda^2\Psi&0&0
\end{array}
\]
In particular $\pi_1(\Gamma^2\Psi)=\Z/2$, whereas
$\pi_1(\Gamma^2\Phi)$, $\pi_1(\Phi\otimes\Psi)$,
$\pi_1(\Psi\otimes\Psi)$ and $\pi_1(\Lambda^2\Psi)$ all vanish.
Applying the additive functor $\pi_1$ to $C$ gives $\Z/2$ in
cochain degree one and zero elsewhere; applying it to $K$ gives
the zero complex. An additive functor preserves a cochain homotopy
and a cochain homotopy equivalence. This proves the asserted
nonexistence of such an equivalence.

For the stronger failure after sections, let $g$ be the nonidentity
element of $\Z/2$ and let $X=\operatorname{sk}_2 B(\Z/2)$.
It has one nondegenerate vertex, one nondegenerate edge $g$, and
one nondegenerate two-simplex $(g,g)$. The outer faces of $(g,g)$
are $g$ and its middle face is degenerate. Thus the reduced
normalized free abelian chains are
\[
 \widetilde N_s\Z[X]
   =\bigl[\Z\xrightarrow{\ 2\ }\Z\bigr]
 \quad\text{in homological degrees }2,1.
\]
For any of the five simplicial modules in the table, a reduced map
$X\to Y$ is exactly a pair $(f_1,f_2)$ with
\[
 f_1\in N_sY_1,\qquad f_2\in N_sY_2,
 \qquad d_0f_1=0,\qquad d_0f_2=2f_1.
\]
This is the exact Dold--Kan correspondence applied to the free
abelian simplicial module $\Z[X]$; no topological replacement of
$X$ is used.

For $Y=\Gamma^2\Phi$ the solution lattice is generated by
$(f_1,f_2)=((1,-2),2)$. For $Y=\Phi\otimes\Psi$ it is
parametrized by
\[
 f_1=(a,b),\qquad f_2=(2a,-2a-2b).
\]
For $\Gamma^2\Psi$ it is parametrized by $(f_1,f_2)=(c,-c)$;
for $\Psi\otimes\Psi$ by
\[
 f_1=e,\qquad f_2=(f,-2e-f);
\]
and for $\Lambda^2\Psi$ by its arbitrary degree-two coordinate.
Substituting these expressions in the vertical differentials gives
exactly the reduced section complexes
\begin{align*}
 \widetilde K(X)&:
 \Z\xrightarrow{\ (1,-2)^{\mathsf t}\ }\Z^2
       \xrightarrow{\ (4,2)\ }\Z,\\
 \widetilde C(X)&:
 \Z\xrightarrow{\ (1,-2,-2)^{\mathsf t}\ }\Z^3
       \xrightarrow{\left(\begin{smallmatrix}0&1&-1\\2&0&1\end{smallmatrix}\right)}
       \Z^2.
\end{align*}
In the first complex the kernel of the last map is generated by
$(1,-2)$ and its image is $2\Z$. Hence its first cohomology
vanishes and its second is $\Z/2$. In the second complex the
kernel of the last map is generated by $(1,-2,-2)$, and the last
two columns of that map have determinant one. It is therefore
surjective, and the whole reduced complex is acyclic. The
basepoint summand contributes only $\Gamma^2(\Z T)$ in cochain
degree zero. This proves the stated difference in second
cohomology.

For the simplicial two-sphere the reduced chain complex is
$\Z[2]$. In $N_s(\Psi\otimes\Psi)_2$, cycles are generated by
$(1,-1)$. The exterior quotient sends that generator to
$2(u_1\wedge u_2)$. The corresponding normalized second cycles
of $\Gamma^2\Phi$, $\Phi\otimes\Psi$ and $\Gamma^2\Psi$
vanish, by the same matrices. Thus both section complexes have
second cohomology $\Z$, and $\theta$ induces multiplication by
two, as claimed.

The independent integer-matrix certificate
\texttt{code/D\_derived\_binomial\_weight2.py} constructs all face
maps from the given bases, applies divided squares, tensor squares
and exterior squares, and checks the proposed normalization and
section lattices as saturated kernels. It also checks the vertical
maps, the Smith forms and the sphere comparison. Its saved output
records $123$ successful assertions. The calculation uses no
floating-point approximation.
\end{proof}

One should also distinguish taking the associated graded before
sections from taking the associated graded after sections. The
functor $\operatorname{Hom}(X,-)$ is left exact and need not
commute with filtration quotients. Statement~\ref{sa08:A022}
settles the explicitly displayed complex $C$ and its underived
sections; it does not silently interchange those two constructions.
